\begin{claim}{B.28}{Proved}{Complete executed extended genus-two catalogue}
The extended mapping class group of the closed orientable genus-two
surface has exactly $92$ conjugacy classes of finite subgroups.  Of
these, $57$ contain the hyperelliptic involution and $35$ do not;
$21$ classes preserve orientation.  The files
\texttt{type\_2\_0\_extended\_catalogue.csv} and
\texttt{type\_2\_0\_extended\_sections.json} specify every group by
an exact finite permutation representation and give every full and
partial splitting, its multiplier and centralizer, all extension
incidences, $\Phi$, and its group orbits and stabilizers.
\end{claim}
\begin{proof}
Let $\iota$ be the hyperelliptic mapping class.  The genus-two
Birman--Hilden homomorphism, with orientation reversal allowed, is
\[
 1\longrightarrow\langle\iota\rangle
 \longrightarrow\mathcal T_{2,0}
 \longrightarrow\mathcal T_{0,6}\longrightarrow1.       \tag{B.28.1}
\]
We recall its geometric content, to specify every comparison used here.
The involution is the deck map of the double cover of $S^2$ branched
at six points.  The covering is determined by the homomorphism from
the fundamental group of the six-times-punctured sphere to $C_2$
that sends each peripheral generator to the nonidentity element;
the product relation is respected because six is even.  Every
homeomorphism preserving the branch set preserves this homomorphism,
and hence has exactly two lifts, differing by $\iota$.
The genus-two hyperelliptic mapping class is central: the usual chain
of five nonseparating curves is invariant under the involution, their
Dehn twists generate the oriented genus-two mapping class group,
and a reflection of the standard branched-cover model commutes with
$\iota$ and supplies the negative coset.  The Birman--Hilden isotopy
property identifies the kernel in (B.28.1) with these two deck classes;
it also makes the lifting map onto.  This is the hyperelliptic
Birman--Hilden description used in B.25, in genus two instead of the
four-point pillowcase; see \cite[\S9.4.2, p.~257, with the proof in \S9.4.3]{bFarbMargalit2012}.

For a finite $K\subset\mathcal T_{2,0}$ put
$E=K\langle\iota\rangle$ and let $H$ be its image in
$\mathcal T_{0,6}$.  The group $E$ is exactly the full inverse image
of $H$ in (B.28.1), and has order $2|H|$.  Thus B.26 supplies its
complete list of $57$ conjugacy classes.  A subgroup $K$ with this
fixed $E$ either equals $E$, or is a complement to
$\langle\iota\rangle$.  Every complement is normal in $E$ and is
uniquely the kernel of a homomorphism
\[
             \chi:E\longrightarrow\mathbf F_2,
             \qquad\chi(\iota)=1.                     \tag{B.28.2}
\]
Indeed $E=K\times\langle\iota\rangle$ for a complement, giving
this homomorphism; conversely the kernel of (B.28.2) has index two
and maps isomorphically onto $H$.  Because $\iota$ is central in the
whole mapping class group, conjugate $K$'s have conjugate associated
$E$'s.  For a fixed $E$, their further fusion is therefore exactly
the action of its normalizer on the characters (B.28.2).

Here are explicit presentations of all the finite groups $E$.
They are also used to calculate the sections, so the statement is
independent of the names assigned to abstract finite groups.  Let
$Q=S^2/H$ be one of the spherical orbifolds in B.26, with its
six-branch-point decoration.  Write $z=\iota$, $z^2=1$, and require
$z$ to commute with all generators.  An interior cone of order $m_j$
has a generator $x_j$ with
\[
                    x_j^{m_j}=z^{d_j},                 \tag{B.28.3}
\]
where $d_j=1$ exactly when its orbit is a branch orbit.  Let $a$ be
the number of regular interior branch orbits.  For a spherical
underlying quotient the remaining relation is
\[
                       \prod_jx_j=z^a.                 \tag{B.28.4}
\]
For a projective-plane quotient, with crosscap generator $u$, it is
$u^2\prod_jx_j=z^a$.

For a mirror disc, expand the cyclic corner list by inserting each
branch orbit on an open mirror edge as an extra entry $(m,d)=(1,1)$.
An original corner contributes $(m,d)$ with its original order $m$
and its branch flag $d$.  The distinction is retained as a color.
Let the resulting cyclic list be $(m_j,d_j)_{j=0}^{L-1}$.
There are reflections $R_0,\ldots,R_{L-1}$ and a boundary element
$e$, with
\[
 \begin{split}
 R_j^2&=1,\\
 \prod_i x_i\,e&=z^a,\\
 (R_jR_{j+1})^{m_j}&=z^{d_j},\\
 R_L&=eR_0e^{-1}.
 \end{split}
 \tag{B.28.5}
\]
When $L=0$ there is one reflection $R_0$ and the last relation is
$eR_0e^{-1}=R_0$.  A newly inserted mirror branch entry says exactly
$R_{j+1}=R_jz$.

These are presentations, rather than only relations satisfied by
some finite images.  To verify this, use the branched cover described
above.  Its full group of lifts of $H$ is a group of order $2|H|$.
A small positive circle about a branch point has deck monodromy $z$;
about an unbranched point it has trivial deck monodromy.  This gives
(B.28.3).  The product of the boundary loops gives (B.28.4), or its
crosscap or disc version.  At a branch point on a mirror, the local
map is $w\mapsto w^2$: the two reflections on the adjacent lifted
half-branches differ by $w\mapsto-w$.  This gives precisely the
inserted entry in (B.28.5).  At a corner the positive local cyclic
generator has its order doubled exactly when it is branched, giving
the other entries.  These local monodromies and $z$ generate the full
lift group.  Conversely quotienting the displayed presentation by
$z=1$ gives the original spherical orbifold presentation, which is
$H$ by the local-injectivity and Euler-characteristic certificate in
B.26.  Since $z$ is central of order at most two, the presented group
has at most $2|H|$ elements.  Its surjection onto the geometric lift
group proves equality and the presentation assertion.

We determine the entire normalizer action needed in (B.28.2), not
just conjugation inside the finite group.  A normalizing mapping
class descends through (B.28.1) and through the $H$-action to an
orbifold homeomorphism of $Q$ preserving its branch decoration.
Conversely any such orbifold homeomorphism lifts to $S^2$, because
$S^2$ is the universal orbifold cover of $Q$, and then to the double
branched cover, because the six peripheral mod-two monodromies are
preserved.  Thus its lifts give the whole normalizer action, with
inner automorphisms and the deck involution as the harmless lifting
ambiguities.  Both act trivially on characters to $\mathbf F_2$.

On a sphere, an orbifold homeomorphism sends each local cone
generator to a conjugate of a local generator with the same original
order and branch flag, or to its inverse when it reverses
orientation.  Characters ignore conjugation and inversion.  All
permutations among cones of the same pair $(m_j,d_j)$ are realized
by homeomorphisms of the sphere.  Point-pushing and pure braids
conjugate the local generators and therefore act trivially on these
characters.  The ordinary branch generators already equal the
central element $z$.  Consequently the exact action on (B.28.2)
is the permutation action on these colored cone coordinates.
The character equations are
\[
       m_j\chi(x_j)=d_j,\qquad
       \sum_j\chi(x_j)=a\pmod2.                       \tag{B.28.6}
\]
This includes, for example, the two elliptic complements in a
four-group lying over a sphere involution with no fixed branch
points: exchanging the two cone points exchanges those complements.
They must be fused, and the program does so.

On a mirror disc there is at most one interior cone.  The normalizer
preserves the cyclic order on its boundary, up to reversal.  It
therefore induces exactly the rotations and reversals of the
expanded colored corner list that preserve its colors.  Each such
symmetry is realized by a disc homeomorphism preserving the interior
decoration.  To see its action explicitly, extend the reflection
sequence by $R_{j+L}=eR_je^{-1}$.  A cyclic shift by $s$ sends
$R_j$ to $R_{j+s}$ and preserves $e$; a reversal sends $R_j$ to
$R_{s-j}$ and $e$ to $e^{-1}$.  The interior cone generator is sent
to a conjugate of itself or its inverse.  Modulo two these formulas
are literal permutations of the reflection coordinates and leave
the interior coordinate fixed.  They respect every relation in
(B.28.5).  Every remaining disc mapping class comes from moving the
interior marked points or from a boundary twist; its effect on local
generators is conjugation and hence is trivial on the characters.
The character equations are therefore
\[
 \begin{split}
 m_i\chi(x_i)&=d_i,\\
 \chi(e)&=a+\sum_i\chi(x_i),\\
 m_j\bigl(\chi(R_j)+\chi(R_{j+1})\bigr)&=d_j\pmod2.
 \end{split}
                                                        \tag{B.28.7}
\]
The program enumerates these bit vectors and divides by exactly the
colored dihedral permutations just proved.  For a projective-plane
quotient, the six-branch-point equation allows only the three cases
$P()$ with $a=3$, $P(2)$ with $(d,a)=(1,1)$, and $P(3)$ with
$(d,a)=(0,1)$.  Their full lifts are respectively $C_4,C_8,C_{12}$,
with $z$ the involution.  None has a character taking $z$ to one.
Thus no further complement fusion is needed in these cases.

Executing (B.28.6)--(B.28.7) gives $35$ character orbits, in addition
to the $57$ full groups.  The following order distribution is a
compact view of the complete $92$-row table:
\[
\begin{array}{c|rrrrrrrr}
 |K|&1&2&3&4&5&6&8&10\\\hline
 \#&1&7&1&13&1&11&14&3\\[2pt]
 |K|&12&16&20&24&32&48&96\\\hline
 \#&17&8&1&9&1&4&1
\end{array}                                             \tag{B.28.8}
\]
In particular its seven involution classes consist of the hyperelliptic
class, the elliptic positive class, and the five negative classes
predicted independently by B.16.

We finish by specifying the verification of all section data.
For each of the $57$ presentations, the program performs an exact
Todd--Coxeter enumeration and obtains the regular action on exactly
$2|H|$ cosets.  It evaluates every relator in these permutations and
checks that $z$ is nonidentity of order two.  In the quotient orbifold
for the genus-two action, an originally branched cone of order $m$
now has order $2m$, a regular branch orbit contributes a cone of
order two, and a branched corner has its order doubled.  Every
local cyclic or dihedral group injects into the computed finite
group.  The orbifold Euler characteristic multiplied by $2|H|$
is $-2$, and the covering subgroup is the kernel of a
monodromy homomorphism compatible with orientation.  Its torsion-free
covering surface is therefore orientable of genus two.
The program also checks that $z$ has exactly six section classes.

The cone and corner cosets, together with the lifted mirror coset
graph of B.26, give every fixed point and reflection oval.  Enumerating
subgroups of these local stabilizers and conjugating them by the
finite group is enough: every finite subgroup with a nonempty fixed
set is contained in such a local stabilizer.  The section multiplier
is read from the power of the corresponding positive local generator;
conjugating by a negative element reverses that angle.  For each
complement, retain exactly the section classes whose image subgroup
lies in its character kernel.  This restriction is exact because
the lift extension for the complement is the inverse image of that
subgroup in the full lift extension.  Recompute maximality in the
retained poset, intersect the section stabilizers with the complement,
and compute its actual orbits.  Thus no maximality or orbit count is
inferred merely by restricting the full-group summary.
All order and power identities for $\mu$ and all unique-maximal-extension
claims for sections containing a rotation are checked.  Finally the
$21$ positive rows agree with the independent primary-source
oriented genus-two catalogue used in B.33 on the complete multiset
of five invariants: group order, full sections, nontrivial partial
sections, $|\Phi|$, and $|\Phi/K|$.
\end{proof}
